\begin{answers}{ch:transformer}{Transformer：ミニGPT を組み立てる}
\item $\mathrm{LN}(\vu)=(-\sqrt2,\,0,\,0,\,\sqrt2)\approx(-1.414,\,0,\,0,\,1.414)$ です。
式\eqref{eq:tf-layernorm}に沿って計算します。
平均は $\mu=(2+4+4+6)/4=4$、平均との差は $(-2,0,0,2)$ です。
分散は $\sigma^2=(4+0+0+4)/4=2$、標準偏差は $\sqrt2\approx1.414$ です。
差を $\sqrt2$ で割ると上の値になります。
すべての成分に 100 を足した $(102,104,104,106)$ でも、結果は\textbf{まったく同じ}です。
平均が 104 に増えるだけで、平均との差 $(-2,0,0,2)$ も分散も変わらないからです。
LayerNorm は、全体を 10 倍する変化（本文の例）だけでなく、全体をずらす変化も消します。

\item 合計は $124{,}439{,}808$（約 1.24 億）で、$12Ld^2$ はその約 68\%（約 7 割）です。
式\eqref{eq:tf-params}の各項は、トークン埋め込み $Vd=50{,}257\times768=38{,}597{,}376$、
位置埋め込み $T_{\max}d=1024\times768=786{,}432$ です。
ブロックは $L(12d^2+4d)=12\times(7{,}077{,}888+3{,}072)=84{,}971{,}520$、
最後の LayerNorm は $2d=1{,}536$、バイアスは $L\times 9d=12\times6{,}912=82{,}944$ です。
足すと $124{,}439{,}808$ になります。
$12Ld^2=84{,}934{,}656$ なので、割合は $84{,}934{,}656/124{,}439{,}808\approx0.683$ です。
残りのほとんど（約 31\%）はトークン埋め込みです。

\item Post-LN の版は、最後の 50 ステップの平均損失が 5.318 で、そのまま（Pre-LN）の 4.165 よりずっと悪くなりました。
本文と同じ設定（8 層、$d=64$、600 ステップ、学習率 $10^{-3}$ で一定、ウォームアップなし）で動かしました。
ブロックの \code{forward} を \code{x = ln1(x + attn(x))}、\code{x = ln2(x + mlp(x))} に変え、最後の LayerNorm は残しました。
50 ステップまでは両者ほぼ同じですが（5.40 と 5.43）、Post-LN は 100 ステップの 5.34 から先ほとんど下がらず、5.3 前後で止まりました。
これは同じデータのユニグラム 5.352 とほぼ同じで、文脈をほとんど使えていないことになります。
\secref{sec:tf-residual}で見たとおり、Post-LN では残差の道の途中に LayerNorm が挟まり、近道が途切れるからだと考えられます。
ウォームアップを入れれば改善するはずですが、ここでは確かめていません。
\item 温度 0 では 20 回すべて「57」で、「47」は 0 回でした。温度 0.8 では「47」が 7 回（残り 13 回は「57」）、温度 1.5 では 9 回でした。
基盤モデルで、乱数の種を 0 に固定して数えました。
「\texttt{=}」の直後の確率は、温度 1 で「5」が 0.559、「4」が 0.428 です。
上位 2 つが接近しているので、貪欲法は毎回わずかに上の「5」を選び、間違えます。
「4」の後に「7」が来る確率は 0.97（温度 0.8）と高いので、「47」の割合はほぼ「4」を引く確率（約 4 割）で決まります。
温度を 1.5 に上げても「4」の確率はほとんど変わりません（0.432）。
しかし\figref{fig:tf-temp}のように裾が持ち上がり、「4」の後の「7」が 0.67 に下がります。
小説の文が続く失敗も 20 回中 4 回出ました。
20 回では回数のばらつきが大きく、7 回と 9 回の差には意味がありません。

\item 困ることは 2 つあります。
1 つ目は因果マスクがないことです。
BERT では各位置が前後のすべてを見て表現を作ります。
1 トークン足すたびに、それまでの全位置の表現が変わるので、毎回すべてを計算し直す必要があります。
2 つ目は学習の目的です。
BERT は、隠した 15\% のトークンを前後から当てる穴埋めで学びました。
次トークン予測のように「$x_{<t}$ だけから $x_t$ の確率」を学んでいないので、条件付き確率を 1 トークンずつ掛け合わせて文全体の確率を作る使い方と合いません。
末尾に穴を置いて繰り返し埋めることはできても、学習と使い方がずれるので、自然な文章になる保証はありません。
\end{answers}
